% !TEX program = lualatex
\documentclass[12pt, a4paper]{article}
\usepackage{master}

\DeclareWorkGR{Cat}{범주론}{Κατηγορίαι}{Categoriae}
\DeclareWorkGR{Met}{형이상학}{Μεταφυσικά}{Metaphysica}
\DeclareWorkGR{Phys}{자연학}{Φυσικά}{Physica}
\DeclareWorkGR{NE}{니코마코스 윤리학}{Ἠθικὰ Νικομάχεια}{Ethica Nicomachea}

\fancyhead[L]{\footnotesize\color{midgray}오르가논 강독}
\fancyhead[R]{\footnotesize\color{midgray}제4주}

\begin{document}
\thispagestyle{firstpage}

\weeksection{제4주}{범주론 6b36--11a38}{관계 범주와 질 범주}

\subsection{관계 범주: 두 정의와 범주의 위기 (7장, 6a36--8a12)}

\subsubsection{관계의 첫 번째 정의}

7장은 범주론에서 가장 많이 논쟁되는 장 가운데 하나입니다.

\srcquote{\gr{Πρός τι δὲ τὰ τοιαῦτα λέγεται, ὅσα αὐτὰ ἅπερ ἐστὶν ἑτέρων εἶναι λέγεται, ἢ ὁπωσοῦν ἄλλως πρὸς ἕτερον.}}{`관계적이라고 불리는 것들은, 자신이 바로 그것인 바가 다른 것의 것이라고 말해지거나, 어떤 다른 방식으로든 다른 것에 대해 말해지는 것들이다.'}{범주론 6a36--37}

예컨대 `두 배'는 어떤 것의 두 배이고, `지식'은 어떤 것에 대한 지식이며, `더 큰'은 어떤 것보다 더 큽니다. 이 정의는 의도적으로 넓습니다 --- `어떤 다른 방식으로든'(\gr{ὁπωσοῦν ἄλλως})이라는 포괄적 조건을 포함합니다.

\subsubsection{관계의 특성들}

아리스토텔레스는 관계에 다음과 같은 특성들이 있다고 분석합니다. 관계 가운데 일부는 반대를 허용합니다 --- 덕과 악덕은 관계적이면서 반대입니다. 다만 두 배와 세 배 같은 관계에는 반대가 없습니다. 관계 가운데 일부는 정도의 차이를 허용합니다 --- `더 비슷한'이라고 말할 수 있습니다.

모든 관계에는 상관자(\gr{ἀντιστρέφον})가 있습니다. 주인이 노예의 주인이라면, 노예는 주인의 노예입니다. 아리스토텔레스는 상관자를 올바르게 짝짓는 것이 중요하다고 강조합니다 --- `날개'를 `새'와 짝지으면 역전이 성립하지 않지만, `날개가 있는 것'과 짝지으면 성립합니다(6b36--7a21).

마지막으로, 관계들은 `본성상 동시적'(\gr{ἅμα τῇ φύσει})입니다 --- 두 배가 있으면 절반이 있고, 절반이 있으면 두 배가 있습니다. 다만 이 원칙에는 예외가 있습니다: 인식 대상(\gr{ἐπιστητόν})은 인식(\gr{ἐπιστήμη})보다 선행합니다(7b22--33).

\subsubsection{범주의 위기: 실체가 관계인가?}

7장의 결정적 전환점은 8a13--28에서 일어납니다. 아리스토텔레스는 첫 번째 정의가 너무 넓어서 실체까지 관계에 포함시킬 위험이 있음을 발견합니다. `머리'는 어떤 것의 머리라고 말해집니다. 그런데 머리는 신체의 부분이고, 신체는 제2실체이므로, 머리는 실체입니다. 그렇다면 첫 번째 정의에 따르면 머리는 관계이기도 합니다.\footnote{이 논변의 재구성은 Matthew Duncombe, ``Aristotle's Two Accounts of Relatives in Categories 7,'' \gri{Phronesis} 60, no.\ 4 (2015): 436--461을 따릅니다.} 이것은 범주의 상호 배타성 원칙에 정면으로 위배됩니다.

\subsubsection{두 번째, 수정된 정의}

이 위기에 대응하여 아리스토텔레스는 관계의 정의를 수정합니다.

\srcquote{\gr{ἔστι δὲ τὰ πρός τι οἷς τὸ εἶναι ταὐτόν ἐστι τῷ πρός τί πως ἔχειν.}}{`관계적인 것들은, 존재함이 어떤 것에 대해 어떤 방식으로 관계를 맺고 있음과 동일한 것들이다.'}{범주론 8a31--32}

`두 배'에게 존재한다는 것은 곧 어떤 것의 두 배이다는 것입니다. 반면 `머리'에게 존재한다는 것은 단순히 어떤 것의 머리라는 것이 아닙니다 --- 머리는 관계 맺음을 넘어서는 고유한 실체적 본성을 가집니다.

{%
\footnotesize
\setlength{\extrarowheight}{1.5pt}%
\renewcommand{\arraystretch}{1.1}%
\setlength{\tabcolsep}{3pt}%
\rowcolors{2}{altrowbg}{white}%
\noindent\centering
\begin{tabular}{|>{\centering\arraybackslash}m{14mm}|>{\centering\arraybackslash}m{30mm}|>{\centering\arraybackslash}m{30mm}|>{\centering\arraybackslash}m{30mm}|}
\hline
\rowcolor{tableheadblue}
{\color{h1navy}\bfseries 항목} &
{\color{h1navy}\bfseries 첫 번째 정의} &
{\color{h1navy}\bfseries 두 번째 정의} &
{\color{h1navy}\bfseries 판정} \\ \hline
두 배 & 통과 & 통과 & 관계 \\ \hline
지식 & 통과 & 통과 & 관계 \\ \hline
머리 & 통과 & 불통과 & 실체 \\ \hline
노예 & 통과 & 통과 & 관계 \\ \hline
\end{tabular}
\par
}

마리오 미뉴치(Mignucci)는 이 수정을 `존재론적 환원 검사'로 해석합니다.\footnote{Mario Mignucci, ``Aristotle's Definitions of Relatives in Cat.\ 7,'' \gri{Phronesis} 31 (1986): 101--127.} 데이비드 세들리(Sedley)는 `경성 관계'(hard relativity)와 `연성 관계'(soft relativity)의 구별로 설명합니다 --- 두 배/절반은 순수하게 상관자와의 관계로 구성되는 경성 관계이고, 머리/신체는 관계를 넘어서는 내재적 조건을 포함하는 연성 관계입니다.\footnote{David Sedley, ``Aristotelian Relativities,'' in \gri{Le Style de la Pensée}, ed.\ M.\ Canto-Sperber and P.\ Pellegrin (Paris: Les Belles Lettres, 2002), 324--352.}

스터트만(Studtmann)이 스탠퍼드 철학 백과사전에서 강조하듯이, 아리스토텔레스의 범주는 현대적 의미의 `관계'(relations)가 아니라 `무엇에 대한 것들'(\gr{τὰ πρός τι})입니다. 안나 마르모도로(Anna Marmodoro)는 이것을 `방향을 가진 단항 속성'(directed monadic properties)이라고 해석합니다 --- 현대적 관계는 복수의 항 사이에 성립하는 다항 속성이지만, 아리스토텔레스의 \gr{πρός τι}는 하나의 항이 다른 항을 향하는 단항 속성이라는 것입니다.\footnote{Anna Marmodoro, \gri{Aristotle on Perceiving Objects} (Oxford: Oxford University Press, 2014).}

\subsection{질 범주: 네 종류와 통일성의 문제 (8장, 8b25--11a38)}

\subsubsection{질의 정의와 네 가지 종류}

\srcquote{\gr{Ποιότητα δὲ λέγω καθ᾽ ἣν ποιοί τινες λέγονται. ἔστι δὲ ἡ ποιότης τῶν πλεοναχῶς λεγομένων. ἓν μὲν οὖν εἶδος ποιότητος ἕξις καὶ διάθεσις λεγέσθω.}}{`질이란 그것에 의해 어떤 성질의 것들이라고 불리는 것을 말한다. 질은 여러 가지 방식으로 말해지는 것 가운데 하나이다. 질의 한 종류는 상태와 상태됨이라 하자.'}{범주론 8b25--27}

\par\nobreak\vspace{-\parskip}\nobreak
{%
\footnotesize
\setlength{\extrarowheight}{1.5pt}%
\renewcommand{\arraystretch}{1.1}%
\setlength{\tabcolsep}{3pt}%
\rowcolors{2}{altrowbg}{white}%
\noindent\centering
\begin{tabular}{|>{\centering\arraybackslash}m{6mm}|>{\centering\arraybackslash}m{32mm}|>{\centering\arraybackslash}m{24mm}|>{\centering\arraybackslash}m{32mm}|}
\hline
\rowcolor{tableheadblue}
{\color{h1navy}\bfseries} &
{\color{h1navy}\bfseries 종류} &
{\color{h1navy}\bfseries 그리스어} &
{\color{h1navy}\bfseries 예시} \\ \hline
1 & 상태와 상태됨 & \gr{ἕξις}/\gr{διάθεσις} & 지식, 덕, 열, 감기 \\ \hline
2 & 자연적 역량·무능력 & \gr{δύναμις}/\gr{ἀδυναμία} & 타고난 권투 선수, 병약함 \\ \hline
3 & 감각적 질·반응 & \gr{παθητικαὶ ποιότητες} & 달콤함, 흼, 열 \\ \hline
4 & 모양과 형태 & \gr{σχῆμα}/\gr{μορφή} & 삼각형, 직선 \\ \hline
\end{tabular}
\par
}

첫째 종류 안에서 상태(\gr{ἕξις})와 상태됨(\gr{διάθεσις})의 차이는 안정성에 있습니다.

\srcquote{\gr{διαφέρει δὲ ἕξις διαθέσεως τῷ μονιμώτερον καὶ πολυχρονιώτερον εἶναι.}}{`상태가 상태됨과 다른 점은 더 안정적이고 더 오래 지속된다는 것이다.'}{범주론 9a14--16}

지식이나 덕처럼 한 번 획득하면 쉽게 잃지 않는 것이 상태(헥시스)이고, 열이나 감기처럼 쉽게 변하는 것이 상태됨(디아테시스)입니다.

\subsubsection{네 종류의 통일성 문제}

이 네 종류가 왜 정확히 이 네 가지인지를 아리스토텔레스는 설명하지 않습니다. 애크릴(Ackrill)은 직설적으로 지적합니다: `그는 상태와 상태됨이 질이라는 특별한 논변을 제시하지 않는다.'\footnote{Ackrill, \gri{Aristotle's Categories and De Interpretatione}, 104.} 몽고메리 퍼스(Furth)는 더 신랄하게 질 범주를 `기괴한 오합지졸'이라고 표현했습니다.\footnote{Montgomery Furth, \gri{Substance, Form and Psyche: An Aristotelean Metaphysics} (Cambridge: Cambridge University Press, 1988).}

스터트만(Studtmann)은 이 비판에 대응하여 질의 네 종류를 두 가지 1차 종류 --- 모양(\gr{σχῆμα})과 성향(디스포지션) --- 로 재편성하는 `재규격화된 해석'을 제안했습니다.\footnote{Paul Studtmann, ``Aristotle's Category of Quality: A Regimented Interpretation,'' \gri{Apeiron} 36, no.\ 3 (2003): 205--227.}

\subsubsection{질의 특성들과 파생어 관계}

질에는 반대가 있을 수 있습니다 --- 정의와 부정의는 반대입니다. 다만 모든 질에 반대가 있는 것은 아닙니다. 질은 정도의 차이를 허용합니다 --- 하나가 다른 것보다 `더 흴' 수 있습니다. 다만 삼각형이 다른 삼각형보다 `더 삼각형'이지는 않습니다(10b26--11a14).

질에 가장 고유한 특성은 비슷함(\gr{ὅμοιον})과 비슷하지 않음(\gr{ἀνόμοιον})입니다(11a15--19). 이것은 3주차에서 다룬 양의 가장 고유한 특성(같음/같지 않음)과 정확히 대응합니다.

아리스토텔레스는 질에서 파생어(\gr{παρώνυμα}) 관계가 작동함을 지적합니다(10a27--b11). 흼(\gr{λευκότης})에서 흰(\gr{λευκόν})이, 문법(\gr{γραμματική})에서 문법적인(\gr{γραμματικός})이 파생됩니다. 이것은 1장에서 도입한 파생어 개념이 실제 범주 분석에서 어떻게 작동하는지를 보여주는 대목입니다.\footnote{M.\ V.\ Wedin, \gri{Aristotle's Theory of Substance: The Categories and Metaphysics Zeta} (Oxford: Oxford University Press, 2000).}

\subsubsection{범주 간 교차의 문제: 질인가 관계인가}

8장의 마지막에서 아리스토텔레스는 7장에서 이미 불거진 문제를 다시 마주합니다(11a20--36). 지식은 질 범주에 속하는 것처럼 보입니다 --- 지식은 영혼의 상태(\gr{ἕξις})이니까요. 그런데 동시에 지식은 관계 범주에도 속하는 것처럼 보입니다 --- 지식은 항상 무엇에 대한 지식이니까요. 아리스토텔레스의 해결책은 이것입니다: 류(\gr{γένος})인 `지식'은 관계적으로 말해질 수 있지만, 종(\gr{εἶδος})인 `문법 지식'은 관계적이 아니라 질적입니다.

이 해결책이 만족스러운지는 논쟁의 여지가 있습니다. 범주의 상호 배타성 원칙은 아리스토텔레스의 범주 체계 전체를 떠받치는 근본 원칙인데, 지식의 사례는 이 원칙에 지속적인 압력을 가합니다.

% ── 참고문헌 ──
\subsection*{참고문헌}
\begin{references}

\refitem Ackrill, J.\ L. \gri{Aristotle's Categories and De Interpretatione}. Clarendon Aristotle Series. Oxford: Clarendon Press, 1963.

\refitem Duncombe, Matthew. ``Aristotle's Two Accounts of Relatives in Categories 7.'' \gri{Phronesis} 60, no.\ 4 (2015): 436--461.

\refitem Furth, Montgomery. \gri{Substance, Form and Psyche: An Aristotelean Metaphysics}. Cambridge: Cambridge University Press, 1988.

\refitem Marmodoro, Anna. \gri{Aristotle on Perceiving Objects}. Oxford: Oxford University Press, 2014.

\refitem Mignucci, Mario. ``Aristotle's Definitions of Relatives in Cat.\ 7.'' \gri{Phronesis} 31 (1986): 101--127.

\refitem Minio-Paluello, L., ed. \gri{Aristotelis Categoriae et Liber de Interpretatione}. OCT. Oxford: Clarendon Press, 1949; repr.\ with corrections, 1956. Eng.\ trans.: J.\ L.\ Ackrill, \gri{Aristotle's Categories and De Interpretatione}, Clarendon Aristotle Series. Oxford: Clarendon Press, 1963.

\refitem Sedley, David. ``Aristotelian Relativities.'' In \gri{Le Style de la Pensée}, edited by M.\ Canto-Sperber and P.\ Pellegrin, 324--352. Paris: Les Belles Lettres, 2002.

\refitem Studtmann, Paul. ``Aristotle's Categories.'' In \gri{Stanford Encyclopedia of Philosophy}.\newline https://plato.stanford.edu/entries/aristotle-categories/.

\refitem ---. ``Aristotle's Category of Quality: A Regimented Interpretation.'' \gri{Apeiron} 36, no.\ 3 (2003): 205--227.

\refitem Wedin, M.\ V. \gri{Aristotle's Theory of Substance: The Categories and Metaphysics Zeta}. Oxford: Oxford University Press, 2000.

\end{references}

\end{document}
